\documentclass[10pt,a4paper]{article}
\usepackage[utf8]{inputenc}
\usepackage[T1]{fontenc}
\usepackage[french]{babel}
\usepackage{lmodern}
\usepackage[left=1.5cm,right=1.5cm,top=1.25cm,bottom=1.25cm]{geometry}
\usepackage{xcolor}
\usepackage{hyperref}
\usepackage{enumitem}
\usepackage{parskip}
\usepackage{fontawesome5}
\usepackage{titlesec}
\usepackage{fancyhdr}
\usepackage{lastpage}

% Couleurs Angular & Python (Indigo profond & Slate)
\definecolor{primary}{RGB}{67, 56, 202}       % Indigo profond (#4338ca)
\definecolor{darkgray}{RGB}{30, 41, 59}       % Texte principal (#1e293b)
\definecolor{slateGray}{RGB}{71, 85, 105}     % Texte secondaire (#475569)
\definecolor{ruleGray}{RGB}{226, 232, 240}    % Ligne séparatrice

\hypersetup{
    colorlinks=true,
    linkcolor=primary,
    urlcolor=primary,
    pdfauthor={Astrid-Aurélien NKUMBE ENONGENE},
    pdftitle={CV Senior Fullstack Angular et Python --- Astrid-Aurélien NKUMBE}
}

% Format des titres
\titleformat{\section}
  {\large\bfseries\color{primary}}
  {}{0em}
  {}[\vspace{-0.35em}\textcolor{primary}{\rule{\linewidth}{1.2pt}}\vspace{-0.25em}]

\titlespacing*{\section}{0pt}{0.85em}{0.4em}

% En-tête et pied de page
\pagestyle{fancy}
\fancyhf{}
\renewcommand{\headrulewidth}{0pt}
\renewcommand{\footrulewidth}{0.4pt}
\renewcommand{\footrule}{\hbox to\headwidth{\color{ruleGray}\leaders\hrule height \footrulewidth\hfill}}
\fancyfoot[L]{\footnotesize\color{slateGray}Astrid-Aurélien NKUMBE --- CV Spécialiste Angular \& Python}
\fancyfoot[C]{}
\fancyfoot[R]{\footnotesize\color{slateGray}Page \thepage\ / \pageref{LastPage}}

\color{darkgray}

\begin{document}

% =========================================================================
% EN-TÊTE PRINCIPAL
% =========================================================================
\begin{center}
    {\LARGE\bfseries\color{primary} Astrid-Aurélien NKUMBE ENONGENE}\\[0.25em]
    {\large\bfseries\color{darkgray} Senior Fullstack Engineer $\cdot$ Angular / TypeScript $\cdot$ Python / FastAPI}\\[0.4em]
    {\footnotesize\color{slateGray}
        \faMapMarker*\ Région Parisienne (Cergy / Paris) --- Mobile France \& Remote \quad$\cdot$\quad
        \faPhone*\ +33 7 44 58 45 62 \quad$\cdot$\quad
        \faEnvelope\ \href{mailto:nkumbeaurelien@hotmail.com}{\texttt{nkumbeaurelien@hotmail.com}}\\[0.25em]
        \faGlobe\ \href{https://nkaurelien.kamitbrains.fr}{\texttt{nkaurelien.kamitbrains.fr}} \quad$\cdot$\quad
        \faGithub\ \href{https://github.com/nkaurelien}{\texttt{github.com/nkaurelien}} \quad$\cdot$\quad
        \faLinkedin\ \href{https://www.linkedin.com/in/nkaurelien/}{\texttt{linkedin.com/in/nkaurelien}}
    }
\end{center}

\vspace{0.1em}
\textcolor{ruleGray}{\rule{\linewidth}{1pt}}
\vspace{0.35em}

% =========================================================================
% PROFIL
% =========================================================================
\section{\faLightbulb\ Profil Professionnel}
Ingénieur Développeur Fullstack Senior avec \textbf{plus de 7 ans d'expérience}, expert sur l'alliance de la robustesse frontend avec \textbf{Angular (TypeScript, RxJS, Signals)} et de la performance backend / data avec \textbf{Python (FastAPI, AsyncIO, architectures RAG/IA)}. Spécialiste de la conception d'architectures applicatives orientées services, du traitement temps réel (WebSockets), de l'interopérabilité des APIs et de l'intégration de modèles d'intelligence artificielle en production.

% =========================================================================
% COMPÉTENCES TECHNIQUES
% =========================================================================
\section{\faCogs\ Compétences Angular \& Python}
\begin{itemize}[leftmargin=1.4em, itemsep=0.25em, topsep=0.15em]
    \item \textbf{Angular \& Écosystème Frontend :} Angular (v12 à v18), TypeScript, RxJS, NgRx, Angular Signals, Angular CLI, Design Systems, HTML5, SCSS.
    \item \textbf{Python \& Frameworks Backend :} Python 3.10+, FastAPI, Pydantic, SQLAlchemy, Alembic, Uvicorn, Gunicorn, AsyncIO, PySpark.
    \item \textbf{IA \& Écosystème Data :} RAG (Retrieval-Augmented Generation), LangChain, LiteLLM, OpenAI API, Ollama, Vector DBs (pgvector), ELK Stack.
    \item \textbf{APIs, Protocoles \& Sécurité :} RESTful APIs, WebSockets, OpenAPI / Swagger, JSON-RPC, normes médicales FHIR/HL7, JWT, OAuth2, Auth0.
    \item \textbf{Bases de Données :} PostgreSQL, MySQL, CouchDB, MongoDB, Elasticsearch, Redis, Supabase.
    \item \textbf{DevOps \& Outillage :} Docker, Docker Compose, Docker Swarm, GitLab CI/CD, Azure DevOps, Nginx, Pytest, Git, ESLint.
\end{itemize}

% =========================================================================
% EXPÉRIENCES PROFESSIONNELLES
% =========================================================================
\section{\faBriefcase\ Expériences Angular \& Python}

\textbf{\large DATA2INNOV} \hfill {\bfseries\small\color{slateGray} Mai 2025 -- Présent} \\
\emph{\color{primary}Lead Développeur Python / FastAPI \& Architecte Données} \hfill {\small France}
\begin{itemize}[leftmargin=1.4em, itemsep=0.18em, topsep=0.15em]
    \item Conception et développement de l'architecture backend asynchrone sous \textbf{FastAPI} et Python 3.11.
    \item Normalisation et validation stricte des flux biomédicaux au format standardisé FHIR / HL7 avec Pydantic.
    \item Ingestion haute fréquence et diffusion temps réel via WebSockets, CouchDB et pipelines Logstash.
    \item Conteneurisation de la stack sous Docker Swarm et supervision de métrologie avec SigNoz.
    \item Mise à jour du support et auteur des fiches pratiques (TP) et fichiers d'évaluation --- Formation \textbf{Cloud Computing \& Conteneurisation}.
\end{itemize}

\vspace{0.35em}

\textbf{\large YSO Conseils / KAZAC} \hfill {\bfseries\small\color{slateGray} Octobre 2023 -- Septembre 2024} \\
\emph{\color{primary}Développeur Python \& Architecte IA (SmartDataPay)} \hfill {\small France}
\begin{itemize}[leftmargin=1.4em, itemsep=0.18em, topsep=0.15em]
    \item Conception du moteur RAG en Python avec LangChain, LiteLLM, OpenAI API et Ollama pour interroger plus de 500 documents complexes.
    \item Création de scripts ETL asynchrones en Python pour l'extraction, le chunking et la vectorisation vers des bases vectorielles PostgreSQL / pgvector.
    \item Développement d'APIs REST modulaires pour l'intégration sécurisée avec les applications web clientes.
\end{itemize}

\newpage

\textbf{\large SYLEN-IT SERVICES} \hfill {\bfseries\small\color{slateGray} Novembre 2022 -- Juillet 2023} \\
\emph{\color{primary}Développeur Angular / TypeScript (Indépendant)} \hfill {\small France}
\begin{itemize}[leftmargin=1.4em, itemsep=0.18em, topsep=0.15em]
    \item Refonte complète des interfaces et de l'architecture composants avec \textbf{Angular}, TypeScript et SCSS pour l'application Coface (groupe BPCE).
    \item Gestion rigoureuse des états applicatifs et flux réactifs asynchrones avec la librairie \textbf{RxJS}.
    \item Intégration de l'API REST Elasticsearch pour l'indexation et la recherche instantanée multicritère de documents.
    \item Migrations de versions Angular, renforcement du typage TypeScript et tests d'intégration.
\end{itemize}

\vspace{0.35em}

\textbf{\large Korée Africa} \hfill {\bfseries\small\color{slateGray} Décembre 2021 -- Juin 2022} \\
\emph{\color{primary}Développeur Backend \& APIs REST} \hfill {\small Remote}
\begin{itemize}[leftmargin=1.4em, itemsep=0.18em, topsep=0.15em]
    \item Développement et sécurisation d'APIs RESTful avec authentification JWT et validation stricte de schémas de données.
    \item Intégration de services Cloud AWS (SQS, SNS, S3) et passerelles de paiement (Stripe, Mobile Money).
    \item Optimisation des temps de réponse grâce à l'indexation Elasticsearch et la mise en cache Redis.
\end{itemize}

\vspace{0.35em}

\textbf{\large MCS IT} \hfill {\bfseries\small\color{slateGray} Juin 2018 -- Octobre 2020} \\
\emph{\color{primary}Développeur Full Stack Angular} \hfill {\small Remote}
\begin{itemize}[leftmargin=1.4em, itemsep=0.18em, topsep=0.15em]
    \item Développement de la plateforme e-commerce Frontend réactive sous \textbf{Angular} et Bootstrap.
    \item Intégration d'APIs REST pour la synchronisation des catalogues produits et le suivi des commandes en temps réel.
\end{itemize}

% =========================================================================
% PROJETS OPEN SOURCE
% =========================================================================
\section{\faRocket\ Projets Open Source \& Innovations}
\begin{itemize}[leftmargin=1.4em, itemsep=0.25em, topsep=0.15em]
    \item \textbf{ai-llm-python-playground (Python, FastAPI) :} Expérimentations d'agents IA, architectures RAG et pipelines LLM modulaires.
    \item \textbf{couchdb-manager (Python, Streamlit) :} Dashboard de supervision, synchronisation et maintenance d'instances CouchDB.
    \item \textbf{fhir-angular-starter (Angular, TypeScript) :} Bibliothèque de services et composants réactifs Angular pour l'intégration des standards FHIR.
\end{itemize}

% =========================================================================
% FORMATIONS & DIPLÔMES
% =========================================================================
\section{\faGraduationCap\ Formations \& Diplômes}
\begin{itemize}[leftmargin=1.4em, itemsep=0.25em, topsep=0.15em]
    \item \textbf{2022 -- 2024} : \textbf{Master of Science (MSc) IA \& Management} / Expert en Ingénierie Informatique (RNCP 7) --- Ionis-STM France.
    \item \textbf{2014 -- 2017} : \textbf{Ingénieur des Travaux en Analyse \& Conception SI} --- Institut Africain de l'Informatique (IAI).
\end{itemize}

\vspace{0.8em}
\begin{center}
    \footnotesize\color{slateGray}
    Retrouvez tous les projets et démos interactives sur \href{https://nkaurelien.kamitbrains.fr}{\textbf{nkaurelien.kamitbrains.fr}}
\end{center}

\end{document}
